% Plantilla de presentación Beamer 16:9 en español.
%   latexmk -pdf -interaction=nonstopmode -halt-on-error presentacion.tex
\documentclass[aspectratio=169,11pt]{beamer}
\usepackage[T1]{fontenc}
\usepackage[spanish]{babel}
\usepackage{lmodern}
\usepackage{booktabs}
\usepackage{siunitx}
\sisetup{output-decimal-marker={,}}
\usetheme{Madrid}
\usecolortheme{beaver}
\setbeamertemplate{navigation symbols}{}
\setbeamertemplate{caption}[numbered]

\title[Título corto]{Título completo del trabajo}
\subtitle{Trabajo Fin de Máster}
\author[N. Apellido]{Nombre Apellido\\ \small Director/a: Nombre Apellido}
\institute[Universidad]{Escuela / Departamento\\ Universidad}
\date{\today}

\begin{document}

\begin{frame}
  \titlepage
\end{frame}

\begin{frame}{Índice}
  \tableofcontents
\end{frame}

\section{Motivación}
\begin{frame}{Por qué importa el problema}
  \begin{itemize}
    \item Contexto en una línea
    \item El problema concreto y su coste
    \item<2-> Lo que falta en el estado del arte
  \end{itemize}
\end{frame}

\section{Objetivos}
\begin{frame}{Objetivos}
  \begin{block}{Objetivo general}
    Desarrollar y validar \dots
  \end{block}
  \begin{enumerate}
    \item Objetivo específico 1
    \item Objetivo específico 2
  \end{enumerate}
\end{frame}

\section{Método}
\begin{frame}{Método en tres pasos}
  \begin{columns}[T]
    \begin{column}{0.5\textwidth}
      \begin{enumerate}
        \item Adquisición de datos
        \item Modelo: $y = \beta_0 + \beta_1 x + \varepsilon$
        \item Validación cruzada ($k=5$)
      \end{enumerate}
    \end{column}
    \begin{column}{0.45\textwidth}
      \centering
      % \includegraphics[width=\linewidth]{figuras/esquema.pdf}
      \fbox{\rule{0pt}{3cm}\rule{0.8\linewidth}{0pt}}
    \end{column}
  \end{columns}
\end{frame}

\section{Resultados}
\begin{frame}{El modelo reduce el error un 35\,\%}
  \centering
  \begin{tabular}{lcc}
    \toprule
    Modelo & RMSE & $R^2$ \\
    \midrule
    Referencia & \num{0.52} & \num{0.71} \\
    Propuesto  & \num{0.34} & \num{0.86} \\
    \bottomrule
  \end{tabular}
\end{frame}

\section{Conclusiones}
\begin{frame}{Conclusiones}
  \begin{enumerate}
    \item Conclusión que responde al objetivo 1
    \item Conclusión que responde al objetivo 2
  \end{enumerate}
  \vfill
  \begin{alertblock}{Trabajo futuro}
    Validación con más datos.
  \end{alertblock}
\end{frame}

\begin{frame}
  \centering\Huge ¡Gracias!\\[1em]
  \normalsize ¿Preguntas?
\end{frame}

\appendix
\begin{frame}{Diapositiva de reserva}
  Detalles para preguntas previsibles del tribunal.
\end{frame}

\end{document}
